\subsection{正函数的积分}

回忆我们已约定把乘积 \(0 \cdot (+\infty)\) 定义为等于 0。这一约定特别有下列推论：若 f 是数值函数 \(\geqslant 0\) （定义在赋有正测度 \(\lambda\) 的局部紧空间上），则 \(\lambda^{*}(af) = a \cdot \lambda^{*}(f)\) 对每个满足 \(0 \leqslant a \leqslant +\infty\) 的常数 a 成立。当 a = 0 时这是明显的；若 \(a = +\infty\) ，则 \(\lambda^{*}(af) = a \cdot \lambda^{*}(f) = 0\) 或 \(\lambda^{*}(af) = a \cdot \lambda^{*}(f) = +\infty\) （视 f 是否为 \(\lambda\) -可忽略而定）；最后，若 \(0 < a < +\infty\) ，则已知 \(\lambda^{*}(af) = a \cdot \lambda^{*}(f)\) 。

命题 5. --- 设 f 是数值函数 \(\geqslant 0\) （在 \(T \times T'\) 上下半连续）。则函数

\[
t \mapsto \int^ {*} f (t, t ^ {\prime}) d \mu^ {\prime} (t ^ {\prime})
\]

在 T 上是下半连续的，且

\[
\iint^ {*} f (t, t ^ {\prime}) d \mu (t) d \mu^ {\prime} (t ^ {\prime}) = \int^ {*} d \mu (t) \int^ {*} f (t, t ^ {\prime}) d \mu^ {\prime} (t ^ {\prime}).\tag{3}
\]

这是 §3, No.~1 命题 2 的推论，其中用到 No.~1 的引理 1。

推论 1. --- 设 \(f\) （相应地，\(f'\) ）是下半连续函数 \(\geqslant 0\) （定义在 \(T\) （相应地，\(T'\) ）上）；则函数 \(f \otimes f': (t, t') \mapsto f(t)f'(t')\) 在 \(T \times T'\) 上下半连续，且

\[
\iint^ {*} f (t) f ^ {\prime} \left(t ^ {\prime}\right) d \mu (t) d \mu^ {\prime} \left(t ^ {\prime}\right) = \left(\int^ {*} f (t) d \mu (t)\right) \left(\int^ {*} f ^ {\prime} \left(t ^ {\prime}\right) d \mu^ {\prime} \left(t ^ {\prime}\right)\right).\tag{6}
\]

用 \(G\) （相应地，\(G'\) ）表示函数 \(g \in \mathcal{H}_{+}(T)\) （相应地，\(g' \in \mathcal{H}_{+}(T')\) ）——其中 \(g \leqslant f\) （相应地，\(g' \leqslant f'\) ）——所成之集；则

\[
f \otimes f ^ {\prime} = \sup _ {g \in \mathrm{G}, g ^ {\prime} \in \mathrm{G} ^ {\prime}} g \otimes g ^ {\prime}.
\]

由于诸函数 \(g \otimes g'\) 属于 \(\mathcal{H}_{+}(\mathrm{T} \times \mathrm{T}')\) ，\(f \otimes f'\) 确实是下半连续的，且 (6) 由命题 5 立即得出（或直接由在前一公式中取极限得出）。

推论 2. --- 设 A 是 T 的一个 \(\mu\) -适度子集，A′ 是 T′ 的一个 \(\mu'\) -适度子集；则 A × A′ 在 T × T′ 中是 \(\nu\) -适度的。

鉴于适度集的定义（§1, No.~2, 命题 5），只须证明：若 B 是 T 中的可积开集，B′ 是 T′ 中的可积开集，则开集 B × B′ 是可积的。这由推论 1 立即得出。

推论 3. --- 设 A 是 T 的一个 \(\mu\) -可忽略子集，B′ 是 T′ 的一个 \(\mu'\) -适度子集；则 A × B′ 是 \(\nu\) -可忽略的。

因为 \(A \times B'\) 是局部 \(\nu\) -可忽略的（命题 4）且是 \(\nu\) -适度的（推论 2），故是 \(\nu\) -可忽略的（ \(\S1\) , No.~2, 命题 7 的推论 1）。

推论 2 与推论 3 可以推广到两个复测度的乘积，只需把该陈述应用于它们的绝对值（Ch. III, §4, No.~2, 命题 3）。

命题 6. --- 设 \(f\) 是数值函数 \(\geqslant 0\) （定义在 \(\mathbf{T} \times \mathbf{T}'\) 上）。则

\[
\iint^ {*} f (t, t ^ {\prime}) d \mu (t) d \mu^ {\prime} (t ^ {\prime}) \geqslant \int^ {*} d \mu (t) \int^ {*} f (t, t ^ {\prime}) d \mu^ {\prime} (t ^ {\prime}).
\]

这由 §3, No.~2 的命题 3 并注意到 (2) 得出。

命题 7. --- 设 \(f\) 是 \(\nu\) -可测的正数值函数（定义在 \(\mathbf{T} \times \mathbf{T}'\) 上）。

\begin{enumerate}
\def\labelenumi{\alph{enumi})}
\tightlist
\item
  若 \(f\) 是 \(\nu\) -适度的，则函数 \(t \mapsto \int^{*} f(t, t') d\mu'(t')\) 、\(t' \mapsto \int^{*} f(t, t') d\mu(t)\) 分别关于 \(\mu\) 和 \(\mu'\) 可测且适度，且
\end{enumerate}

\[
\begin{array}{r} \iint^ {*} f (t, t ^ {\prime}) d \mu (t) d \mu^ {\prime} (t ^ {\prime}) = \int^ {*} d \mu (t) \int^ {*} f (t, t ^ {\prime}) d \mu^ {\prime} (t ^ {\prime}) \\ = \int^ {*} d \mu^ {\prime} (t ^ {\prime}) \int^ {*} f (t, t ^ {\prime}) d \mu (t). \end{array}\tag{7}
\]

\begin{enumerate}
\def\labelenumi{\alph{enumi})}
\setcounter{enumi}{1}
\tightlist
\item
  若测度 \(\mu'\) 是适度的，则函数 \(t \mapsto \int^{\bullet} f(t, t') d\mu'(t')\) 是 \(\mu\) -可测的，且
\end{enumerate}

\[
\iint^ {\bullet} f (t, t ^ {\prime}) d \mu (t) d \mu^ {\prime} (t ^ {\prime}) = \int^ {\bullet} d \mu (t) \int^ {\bullet} f (t, t ^ {\prime}) d \mu^ {\prime} (t ^ {\prime}).\tag{8}
\]

断言 a) 以及 \(\mu'\) 有界时的断言 b) 都是 §3, No.~2 命题 5 的推论。为处理 \(\mu'\) 适度的情形，把 \(\mu'\) 表示成有界测度序列之和 \(\sum_{n\in\mathbf{N}}\mu_n'\) （§2, No.~3, 命题 4）。于是函数 \(t\mapsto\int^{\bullet}f(t,t')d\mu_n'(t')\) 是 \(\mu\) -可测的，且

\[
\iint^ {\bullet} f (t, t ^ {\prime}) d \mu (t) d \mu_ {n} ^ {\prime} (t ^ {\prime}) = \int^ {\bullet} d \mu (t) \int^ {\bullet} f (t, t ^ {\prime}) d \mu_ {n} ^ {\prime} (t ^ {\prime}).
\]

但 \(\mu \otimes \mu' = \sum_{n \in \mathbf{N}} (\mu \otimes \mu_n')\) （命题 1）；于是对 \(n\) 求和即得断言 b)（ \(\S 2\) , No.~2, 命题 1）。

推论 1. --- 设 H 是 T × T′ 的一个子集，并设 A 是由使 H 在 t 处的截口 H(t) 不是 μ′-可忽略的 t ∈ T 组成之集。

\begin{enumerate}
\def\labelenumi{\alph{enumi})}
\item
  若 H 是 \(\nu\) -可忽略的，则 A 是 \(\mu\) -可忽略的。
\item
  若 H 是局部 \(\nu\) -可忽略的且 \(\mu'\) 是适度的，则 A 是局部 \(\mu\) -可忽略的。
\end{enumerate}

性质 a) 由命题 7（或命题 6）立即得出。在 b) 的假设下，说 \(H(t)\) 局部 \(\mu'\) -可忽略与说它 \(\mu'\) -可忽略是一回事，因为 \(\mu'\) 是适度的（ \(\S1\) , No.~2, 命题 7）。于是性质 b) 由公式 (8) 得出。

这一推论通过取绝对值立即推广到两个复测度的乘积。对下列推论同样成立：

推论 2. --- 若集 \(A \subset T \times T'\) 是 \(\nu\) -可积的，则 A 在 t 处的截口 \(\mathrm{A}(t)\) 是 \(\mu'\) -可积的（对几乎每个 \(t \in T\) ），函数 \(t \mapsto \mu'(\mathrm{A}(t))\) 是 \(\mu\) -可积的，且

\[
\nu (\mathrm{A}) = \int \mu^ {\prime} (\mathrm{A} (t)) d \mu (t).\tag{9}
\]

命题 8. --- 对每对分别定义在 T 和 T′ 上的数值函数 \(f \geqslant 0\) 、\(f' \geqslant 0\) ，

\[
\iint^ {\bullet} f (t) f ^ {\prime} \left(t ^ {\prime}\right) d \mu (t) d \mu^ {\prime} \left(t ^ {\prime}\right) = \left(\int^ {\bullet} f (t) d \mu (t)\right) \left(\int^ {\bullet} f ^ {\prime} \left(t ^ {\prime}\right) d \mu^ {\prime} \left(t ^ {\prime}\right)\right).\tag{10}
\]

我们先处理 \(\mu\) 和 \(\mu'\) 是具紧支集测度的情形；此时 \(\mu \otimes \mu'\) 亦然，且所有符号 \(\int^{\bullet}\) 、\(\iint^{\bullet}\) 都可换成上积分。由命题 6，

\[
\begin{array}{r l} \iint^ {*} f (t) f ^ {\prime} (t ^ {\prime}) d \mu (t) d \mu^ {\prime} (t ^ {\prime}) & \geqslant \int^ {*} d \mu (t) \int^ {*} f (t) f ^ {\prime} (t ^ {\prime}) d \mu^ {\prime} (t ^ {\prime}) \\ & = \left(\int^ {*} f (t) d \mu (t)\right) \left(\int^ {*} f ^ {\prime} (t ^ {\prime}) d \mu^ {\prime} (t ^ {\prime})\right). \end{array}
\]

为建立反向不等式，选取一个函数 \(h \geqslant f\) （相应地，\(h' \geqslant f'\) ），它是一个下半连续函数序列 \((h_n)\) （相应地，\((h'_n)\) ）的下包络，使得

\[
\int^ {*} h (t) d \mu (t) = \int^ {*} f (t) d \mu (t)
\]

\((\text{resp. } \int^{*} h'(t') d\mu'(t') = \int^{*} f'(t') d\mu'(t'))\) ；这样的函数的存在性由上积分的定义（Ch. IV, §1, No.~3, 定义 3）和 Lebesgue 定理立即得出。把命题 7 应用于可测函数 \(h \otimes h'\) ，有

\[
\begin{array}{l} \iint^ {*} f (t) f ^ {\prime} (t ^ {\prime}) d \mu (t) d \mu^ {\prime} (t ^ {\prime}) \leqslant \iint^ {*} h (t) h ^ {\prime} (t ^ {\prime}) d \mu (t) d \mu^ {\prime} (t ^ {\prime}) \\ = \left(\int^ {*} h (t) d \mu (t)\right) \left(\int^ {*} h ^ {\prime} (t ^ {\prime}) d \mu^ {\prime} (t ^ {\prime})\right) \\ = \left(\int^ {*} f (t) d \mu (t)\right) \left(\int^ {*} f ^ {\prime} (t ^ {\prime}) d \mu^ {\prime} (t ^ {\prime})\right), \end{array}
\]

这就是所要的不等式。于是当 \(\mu\) 和 \(\mu'\) 是具紧支集测度时命题得证。为处理一般情形，只须把 \(\mu\) （相应地，\(\mu'\) ）表示成具紧支集测度的族 \((\mu_{\alpha})_{\alpha\in\mathbb{A}}\) （相应地，\((\mu'_{\beta})_{\beta\in\mathbb{B}}\) ）之和（ \(\S2\) , No.~3, 命题 4），对每个测度 \(\mu_{\alpha}\otimes\mu'_{\beta}\) 写出公式 (10)，再对 \((\alpha,\beta)\) 求和，其中用到命题 1（ \(\S2\) , No.~2, 命题 1）。

推论 1. --- 沿用命题 8 中的记号，

\[
\iint^ {*} f (t) f ^ {\prime} \left(t ^ {\prime}\right) d \mu (t) d \mu^ {\prime} \left(t ^ {\prime}\right) = \left(\int^ {*} f (t) d \mu (t)\right) \left(\int^ {*} f ^ {\prime} \left(t ^ {\prime}\right) d \mu^ {\prime} \left(t ^ {\prime}\right)\right)\tag{11}
\]

但可能除去第二端的一个因子等于 0 而另一个等于 +∞的情形。

当第二端的两个因子都有限时，函数 \(f\) 和 \(f'\) 都是适度的（§1, No.~2, 命题 7），因此函数 \(f \otimes f'\) 是适度的（命题 5 的推论 2）；于是上述等式化为公式 (10)（§1, No.~2, 命题 7）。当第二端的一个因子取值 \(+\infty\) 而另一个不为零时，第二端取值 \(+\infty\) ，上述等式由命题 6 得出。

推论 2. --- 设 f 和 \(f'\) 是两个取值于 C 或 \(\overline{R}\) 中的函数，分别定义在 T 和 \(T'\) 上，且分别关于测度 \(\mu\) 和 \(\mu'\) 本质可积（相应地，可积）。则函数 \(f \otimes f'\) 关于测度 \(\mu \otimes \mu'\) 本质可积（相应地，可积），且

\[
\iint f (t) f ^ {\prime} \left(t ^ {\prime}\right) d \mu (t) d \mu^ {\prime} \left(t ^ {\prime}\right) = \left(\int f (t) d \mu (t)\right) \left(\int f ^ {\prime} \left(t ^ {\prime}\right) d \mu^ {\prime} \left(t ^ {\prime}\right)\right).\tag{12}
\]

当 \(f\) 和 \(f'\) 为正时，\(f \otimes f'\) 依命题 3 的推论 3 可测，且该陈述由公式 (10)（相应地，(11)）以及本质可积性判别法（§1, No.~3, 命题 9）（相应地，Ch. IV, §5, No.~6 的可积性判别法，定理 5）得出。一般情形随后立即得出。

推论 2 立即推广到两个复测度的乘积。

